\section{Portée des résultats, limites et perspectives}
\label{sec:discussion}

\subsection{Réponse conditionnelle à la question de comparaison}

La diffusion dépendant des variations locales modifie effectivement le
compromis entre réduction du bruit et maintien de structures dans la
grille étudiée, mais cette modification dépend du seuil, de la conductance,
du temps et du critère. Au cas fixé, la rationnelle $K=0{,}10$ obtient
de meilleurs PSNR et SSIM que la chaleur à $n=20$. L'exponentielle au
même seuil améliore seulement le premier de ces deux critères.
À faible $K$, la limitation des flux peut aussi conserver le bruit ;
à forte diffusion ou après trop d'itérations, PM comme la chaleur peut
sur-lisser. Le gain potentiel n'est donc pas une propriété inconditionnelle
du nom de la méthode.

Ces observations sont compatibles avec le mécanisme des flux, sans
valoir démonstration d'une théorie de performance. La conservation de
la moyenne et le principe du maximum concernent l'évolution discrète
elle-même ; la fidélité à une image propre donnée relève d'une autre
question. Une évolution qui reste bornée entre les extrema de l'entrée
peut néanmoins éloigner les contours de leur position ou de leur largeur
de référence. De même, la dissipation de la chaleur ne garantit pas une
diminution de l'erreur de reconstruction à chaque étape.

\subsection{Ce que le plan ne permet pas d'affirmer}

Deux images synthétiques mettent en évidence des mécanismes, mais ne
représentent pas la diversité des images naturelles, de leurs textures
et de leurs appareils d'acquisition. Dix réalisations de bruit sur une
image ne remplacent pas dix images. Le bruit est exclusivement gaussien
additif indépendant par pixel ; les bruits dépendant de l'intensité,
impulsionnels ou spatialement corrélés ne sont pas étudiés. La réutilisation
du champ standardisé par graine et la construction du protocole après
la phase exploratoire limitent également l'interprétation statistique.

La grille de $K$ et de temps est finie. Ses maxima ne sont pas des
optima globaux, surtout lorsqu'ils atteignent le dernier checkpoint.
Les oracles exploratoires utilisent la vérité terrain : ils renseignent
le potentiel observé sur cette grille, sans proposer une règle applicable
à une photographie inconnue. Les PSNR et SSIM portent des notions de
fidélité différentes ; aucun ne remplace une validation perceptuelle ni
une analyse géométrique dédiée des contours. Les cartes de gradient
restent un support d'interprétation visuelle.

Enfin, la PM exécutée est directionnelle à quatre voisins. Ses flux ne
reconstruisent pas une norme isotrope complète du gradient aux interfaces,
et l'orientation de la grille peut intervenir. La stabilité démontrée
sur cette grille ne résout pas les difficultés du problème continu
non linéaire. Aucune contraction entre solutions, convergence globale
ou décroissance d'une énergie PM entièrement discrète n'est déduite des
figures et tests réalisés.

\subsection{Perspectives non réalisées}

Une prolongation pertinente serait de valider les observations sur des
images naturelles dont la provenance et les conditions d'utilisation
sont vérifiées. Elle devrait distinguer réglage des paramètres et
évaluation sur de nouvelles images. La régularisation proposée par
Catté et collaborateurs \cite{catte1992}, ainsi que d'autres approches
de diffusion discutées par Weickert \cite{weickert1998}, offrent des
pistes pour traiter les difficultés du modèle continu et la dépendance
à la grille ; elles n'ont pas été implémentées ici.

Une sélection de $K$ et de l'arrêt sans accès à l'image propre est un
problème à part entière. Elle pourrait mobiliser une estimation du bruit,
des critères de discrépance ou une autre stratégie explicitement validée.
Les liens avec problèmes inverses et optimisation permettraient aussi
de distinguer l'évolution de lissage d'une reconstruction guidée par un
modèle d'observation. Des comparaisons avec d'autres schémas numériques
ou des méthodes apprises exigeraient un protocole nouveau et des budgets
comparables. Ces pistes sont des perspectives, non des résultats de ce
projet.
